\documentclass[11pt,a4paper]{article}
\usepackage{amsmath,amssymb}
\usepackage{hyperref}
\usepackage{booktabs}
\usepackage{geometry}
\geometry{margin=2.5cm}

\newcommand{\MeV}{\,\mathrm{MeV}}
\newcommand{\GeV}{\,\mathrm{GeV}}
\newcommand{\TeV}{\,\mathrm{TeV}}

\title{What is the inverse Koide formula?}
\author{A.~Rivero\\
  \small\textit{Universidad de Zaragoza}\\
  \small\texttt{rivero@unizar.es}}
\date{\today}

\begin{document}
\maketitle

\begin{abstract}
We give a pedestrian introduction to the ``inverse Koide formula''
for quark masses, explain why it exists alongside the original
lepton formula, and show that both follow from a single algebraic
mechanism in a confining gauge theory.
Along the way we encounter an operator that maps leptons to quarks,
a potential that stubbornly prefers democracy, and a 65-year-old
theorem by Kostant that selects three families of matter.
\end{abstract}


\section{The original Koide formula (1983)}

In 1983, Yoshio Koide noticed something peculiar about the charged
lepton masses~\cite{Koide:1983}.
Define
\begin{equation}
  K \;=\; \frac{m_e + m_\mu + m_\tau}
  {\bigl(\sqrt{m_e} + \sqrt{m_\mu} + \sqrt{m_\tau}\bigr)^2}\,.
  \label{eq:K}
\end{equation}
Plug in the numbers ($m_e = 0.511\MeV$, $m_\mu = 105.66\MeV$,
$m_\tau = 1776.86\MeV$) and you get
\begin{equation}
  K = 0.666661\ldots
\end{equation}
which is $2/3$ to $0.001\%$.

This is strange. The ratio $K$ can range from $1/3$ (all masses equal)
to $1$ (one mass dominates). There is no obvious reason for it to be
$2/3$.

Koide's formula has survived 40 years of improved measurements.
The $\tau$ mass has been remeasured many times since 1983; each time,
$K$ has moved \emph{closer} to $2/3$, not farther.
The current precision makes accidental coincidence implausible
($p < 10^{-5}$ from Monte Carlo).


\section{What are ``signed charges''?}

The key to understanding $K = 2/3$ is to change variables.
Define $q_i = \sqrt{m_i}$ (the ``signed charge'').
Then
\begin{equation}
  K = \frac{q_1^2 + q_2^2 + q_3^2}{(q_1 + q_2 + q_3)^2}\,.
  \label{eq:Kq}
\end{equation}
This is just a ratio of power sums in the $q_i$.
The numerator is the sum of squares; the denominator is the
square of the sum.

Now the Cauchy--Schwarz inequality gives $1/3 \leq K \leq 1$,
with $K = 1/3$ when all $q_i$ are equal and $K = 1$ when
one dominates.
The value $K = 2/3$ is the midpoint --- the ``equipartition''
between uniformity and hierarchy.


\section{What is the \emph{inverse} Koide formula?}

Take the same ratio~\eqref{eq:Kq} but with $q_i \to 1/q_i$:
\begin{equation}
  K^{-1} \;\equiv\;
  \frac{1/m_d + 1/m_s + 1/m_b}
  {\bigl(1/\sqrt{m_d} + 1/\sqrt{m_s} + 1/\sqrt{m_b}\bigr)^2}\,.
  \label{eq:Kinv}
\end{equation}
Plug in the down-type quark masses ($m_d = 4.67\MeV$,
$m_s = 93.4\MeV$, $m_b = 4183\MeV$) and you get
\begin{equation}
  K^{-1}(d,s,b) = 0.6652\,,
\end{equation}
which is $2/3$ to $0.22\%$.
Not as precise as the leptons, but far better than random
($p \sim 0.3\%$ for a single triple).

Under renormalisation-group running (5-loop Standard Model,
using the Martin--Robertson SMDR code~\cite{MartinRobertson:2019}),
the inverse down-type ratio improves monotonically toward the UV:

\begin{center}
\begin{tabular}{lll}
\toprule
Scale $Q$ & $K^{-1}(d,s,b)$ & Deviation from $2/3$ \\
\midrule
$M_Z$        & $0.66750$ & $+0.13\%$ \\
$1\TeV$      & $0.66719$ & $+0.08\%$ \\
$100\TeV$    & $0.66675$ & $+0.01\%$ \\
$280\TeV$    & $\mathbf{0.66667}$ & $\mathbf{0.00\%}$ \\
$10^6\GeV$   & $0.66657$ & $-0.01\%$ \\
$M_\text{Planck}$ & $0.66539$ & $-0.19\%$ \\
\bottomrule
\end{tabular}
\end{center}
The inverse ratio crosses $2/3$ exactly at $Q = 280\TeV$.
This is not a fit --- it is a consequence of how down-type
Yukawa couplings run in the Standard Model.

There is also a second inverse triple: $K^{-1}(s,c,-t)$
crosses $2/3$ at $Q = 82\TeV$ (the $-$ sign on $t$ means the
top charge enters with opposite sign, which is needed for
the formula to work).


\section{Why should there be a direct \emph{and} an inverse formula?}

In supersymmetric QCD with three colours and three flavours
($\mathrm{SU}(3)$ with $N_f = 3$), the theory
s-confines~\cite{Seiberg:1995}.
The low-energy degrees of freedom are ``mesons''
$M^i{}_j = Q^i \tilde{Q}_j$, subject to the quantum constraint
\begin{equation}
  \det M = \Lambda^6\,.
\end{equation}
The vacuum expectation values of the mesons satisfy a seesaw:
\begin{equation}
  \langle M_i \rangle = \frac{\Lambda^2 P^{1/3}}{m_i}\,,
  \qquad P = m_1 m_2 m_3\,.
\end{equation}
Heavy preons give light mesons, and vice versa.
This is the ``Seiberg seesaw.''

But the \emph{same} theory also contains the adjugate:
\begin{equation}
  \mathrm{adj}(M)_i = \frac{\Lambda^4}{P^{1/3}}\,m_i\,,
\end{equation}
which goes the other way: heavy preons give heavy adjugates.

So one theory produces \textbf{two} operators:
\begin{itemize}
  \item The meson VEV: $\langle M_i \rangle \propto 1/m_i$
    (``inverse'')
  \item The adjugate: $\mathrm{adj}(M)_i \propto m_i$ (``direct'')
\end{itemize}
If the ultraviolet masses satisfy $m_i = \mu\,q_i^2$ and
$K(q_1, q_2, q_3) = 2/3$ (which is automatic for any
$\mathrm{U}(3)$ charges), then \emph{both} the direct and the
inverse ratios equal $2/3$.
The inverse Koide is not a separate coincidence --- it is the
other face of the same coin.

The two operators are linked by a family-independent product rule:
\begin{equation}
  \langle M_i \rangle \cdot \mathrm{adj}(M)_i
  = \Lambda^6 \qquad \text{(same for all } i\text{)}\,.
\end{equation}
This is the constraint $\det M = \Lambda^6$ in disguise.


\section{The operator that maps leptons to quarks}

Define the matrix $R$ by
\begin{equation}
  R_{ij} = \langle M_k \rangle
  \quad (k \neq i, j;\; i \neq j)\,,\qquad
  R_{ii} = 0\,.
\end{equation}
This $3 \times 3$ symmetric, traceless matrix satisfies an
exact identity:
\begin{equation}
  R \cdot \mathbf{M} = 2\,\mathrm{adj}(\mathbf{M})\,.
\end{equation}
In words: $R$ acting on the meson VEV vector gives twice the
adjugate vector.
If you identify mesinos (meson fermions) with leptons and
adjugate fermions with quarks, then $R$ is the operator that
\emph{rotates leptons into quarks}.

What algebra does $R$ satisfy?
The Cayley--Hamilton theorem on the constraint surface gives:
\begin{equation}
  R^3 = \bigl(\textstyle\sum M_i^2\bigr)\,R + 2\Lambda^6\,I\,.
\end{equation}
This is a \emph{cubic}, not the quadratic $J^2 = -1$ of
standard $\mathrm{SU}(2)_R$ in $\mathcal{N} = 2$ supersymmetry.
The ``emergent $\mathcal{N} = 2$'' of the sBootstrap is
therefore not standard $\mathcal{N} = 2$ --- it is a cubic
deformation controlled by the mass hierarchy.

A particularly beautiful property: conjugating by the diagonal
meson matrix $D = \mathrm{diag}(M_1, M_2, M_3)$ gives
\begin{equation}
  D\,R\,D = \Lambda^6\,(J - I)\,,
\end{equation}
where $J$ is the all-ones matrix. The right-hand side is the
\textbf{democratic matrix} --- completely independent of the
individual masses. All hierarchy information is absorbed by
the $D$ conjugation; the $R$ structure itself is universal.


\section{Why three families?}

Here is a curious coincidence (or not).
For $\mathrm{U}(N)$ with canonical normalisation, the Koide ratio
generalises to $K_N = 2/N$.
Meanwhile, the principal $\mathfrak{sl}(2)$
subalgebra~\cite{Kostant:1959} of $\mathfrak{sl}(N)$ has mean-square
weight
\begin{equation}
  \langle J_3^2 \rangle = \frac{N^2 - 1}{12}
\end{equation}
in the fundamental representation.

Setting $K_N = \langle J_3^2 \rangle$:
\begin{equation}
  \frac{2}{N} = \frac{N^2 - 1}{12}
  \qquad\Longleftrightarrow\qquad
  N(N^2 - 1) = 24\,.
\end{equation}
This factors as $(N-3)(N^2 + 3N + 8) = 0$.
The quadratic has no real roots ($\Delta = 9 - 32 < 0$).

The unique positive integer solution is $N = 3$.

In other words: the Koide value $2/3$ is the Casimir eigenvalue
of the principal $\mathfrak{sl}(2)$ in $\mathfrak{sl}(N)$, and
this identification works \emph{only} for three families.
Kostant proved in 1959 that this principal
$\mathfrak{sl}(2)$ is the unique canonical subalgebra of any
simple Lie algebra~\cite{Kostant:1959}.
He did not know about quarks, leptons, or the Koide formula.
The connection was noticed 65 years later.


\section{What we don't know}

Lest the reader think the problem is solved, here is an honest
list of what remains open:

\begin{enumerate}
  \item \textbf{The angle.} The charge parametrisation has one
    free parameter (the ``Cartan angle'' $\alpha \approx 12.7°$)
    that determines all mass ratios within a sector. Nobody knows
    where it comes from. A one-loop Coleman--Weinberg analysis
    shows that every perturbative contribution to the potential
    prefers $\alpha = 0$ (all masses equal). The physical value
    is a flat direction of the perturbative theory.

  \item \textbf{The CKM matrix.} The off-diagonal meson masses
    in the confined theory block-decouple from the diagonal ones
    at tree level. Quark mixing cannot arise from tree-level meson
    mixing; it requires loop effects or K\"ahler corrections.

  \item \textbf{Neutrinos.} The tree-level prediction for the
    PMNS matrix is diagonal (no mixing), which is wrong.

  \item \textbf{The cubic algebra.} The $R$ operator satisfies a
    cubic, not a quadratic. Does this define a new kind of
    extended supersymmetry?

  \item \textbf{Why $m \propto q^2$?} The charge-squared mass
    law is uniquely selected by an algebraic theorem (among
    monomials, only $n = 2$ gives a direction-independent $K$).
    But no one has derived it from the dynamics of the confining
    theory. It is the framework's central postulate.
\end{enumerate}

Five open problems. Always five.


\section*{Summary}

\begin{center}
\renewcommand{\arraystretch}{1.3}
\begin{tabular}{lll}
\toprule
Quantity & Value & Precision \\
\midrule
$K(e, \mu, \tau)$ & $0.666661$ & $0.001\%$ \\
$K^{-1}(d, s, b)$ at $280\TeV$ & $0.666667$ & exact \\
$K^{-1}(s, c, -t)$ at $82\TeV$ & $0.666667$ & exact \\
$m_b$ from bion relation & $4183\MeV$ & $0.01\%$ \\
$m_t$ from chain completion & $168.5\GeV$ & $2.3\%$ \\
$|V_{us}|$ from $\sqrt{m_\mu/m_\tau}$ & $0.244$ & $8.7\%$ \\
\midrule
$N = 3$ from $K = \langle J_3^2 \rangle$ & unique & exact \\
\bottomrule
\end{tabular}
\end{center}

The inverse Koide formula is not a separate observation.
It is the shadow of the same algebraic structure that produces
the original Koide formula, viewed through the Seiberg seesaw
of a confining gauge theory. Both the direct and inverse
relations are automatic for $\mathrm{U}(3)$ family charges.
The value $2/3$ is Kostant's Casimir, and it selects three
families from a Diophantine equation that Bertram Kostant never
had reason to write down.


\begin{thebibliography}{9}

\bibitem{Koide:1983}
  Y.~Koide,
  Phys.\ Lett.\ B \textbf{120}, 161 (1983).

\bibitem{MartinRobertson:2019}
  S.P.~Martin and D.G.~Robertson,
  Phys.\ Rev.\ D \textbf{100}, 073004 (2019),
  arXiv:1907.02500.

\bibitem{Seiberg:1995}
  N.~Seiberg,
  Nucl.\ Phys.\ B \textbf{435}, 129 (1995),
  arXiv:hep-ph/9411149.

\bibitem{Kostant:1959}
  B.~Kostant,
  Amer.\ J.\ Math.\ \textbf{81}, 973 (1959).

\bibitem{Rivero:2024}
  A.~Rivero,
  Eur.\ Phys.\ J.\ C \textbf{84}, 1058 (2024),
  arXiv:2407.05397.

\end{thebibliography}

\end{document}
